\documentclass[12pt]{article}
\usepackage{fontspec}
\usepackage{graphicx}
\usepackage{pst-plot}
\usepackage{scrextend}
\usepackage{advdate}
\usepackage{listings}
\usepackage{tikz}
\usepackage{amssymb}
\usepackage{color}
\usepackage{soul}
\usepackage{caption}
\usepackage{fancyhdr}
\usepackage{pgfornament}
\usepackage[many]{tcolorbox}
\usepackage{collectbox}
\usepackage{mdframed}
\usepackage{xcolor}
\usepackage{mathtools}
\usepackage[hidelinks]{hyperref}
\usepackage[final]{pdfpages}
\usepackage[left=0.5in,right=0.5in,top=0.5in,bottom=0.5in,footskip=15mm]{geometry}

\newfontfamily{\ubuntu}[Path=./fonts/,]{UbuntuMono-Regular}
\newfontfamily{\ubuntui}[Path=./fonts/,]{UbuntuMono-Italic}
\newfontfamily{\ubuntub}[Path=./fonts/,]{UbuntuMono-Bold}

\definecolor{blue1}{RGB}{115, 3, 252}
\definecolor{blue2}{RGB}{3, 53, 252}
\definecolor{orange1}{RGB}{252, 107, 3}

\begin{document}
\pagenumbering{gobble}

\vspace{5mm}


\noindent {\ubuntu \textbf{\textcolor{red}{DUE:}} Monday, May 08, 2023} \hfill {\ubuntu\textbf{\color{blue1}Rocco DeVito}}

\hfill {\ubuntu Issued \AdvanceDate[0]\today}

\vspace{10mm}

\noindent {\Large \textbf{\textcolor{blue1}{{\ubuntu Cyclic Groups}  \hrulefill}}}

\vspace{10mm}

\noindent{\ubuntu A group is {\ubuntub cyclic} if all the elements of the group can be generated by repeatedly applying the group operation to a single element of the group. This element is called a {\ubuntub generator} of the group. Let's look at an example of a cyclic group.}

\vspace{10mm}

\noindent{\ubuntub\textcolor{blue1}{Example: }}{\ubuntu We already know that $\left(\mathbb{Z}/5\mathbb{Z}\right)^{\times} = \left\lbrace 1, 2, 3, 4 \right\rbrace$ is a group under multiplication modulo 5. It also happens to be cyclic because all of its elements can be generated by the element $2$ (in fact, $3$ can also generate the group). Here's what I mean:}

\begin{align*}
    2 &= 2 \\
    2 \times 2 \pmod{5} &= 4 \pmod{5} = 4 \\
    2 \times 2 \times 2 \pmod{5} &= 8 \pmod{5} = 3 \\
    2 \times 2 \times 2 \times 2 \pmod{5} &= 16 \pmod{5} = 1. 
\end{align*}

\vspace{2mm}

\noindent{\ubuntu We can generate every element in $\left\lbrace 1, 2, 3, 4 \right\rbrace$ by multiplying $2$ by itself $\pmod{5}$ over and over again. Because of this, we can write $$\left(\mathbb{Z}/5\mathbb{Z}\right)^{\times} = \left<2\right>.$$}

\noindent{\ubuntu This notation tells us that $\left(\mathbb{Z}/5\mathbb{Z}\right)^{\times}$ is a cyclic group generated by the element $2$.}

\vspace{10mm}

\begin{center}
{\Large \textbf{\textcolor{blue1}{{\ubuntu Exercises}}}}
\end{center}

\vspace{5mm}

\noindent {\ubuntub (1.)} {\ubuntu Verify that $\left(\mathbb{Z}/5\mathbb{Z}\right)^{\times} = \left< 3 \right>$, i.e., that $3$ is also a generator of $\left(\mathbb{Z}/5\mathbb{Z}\right)^{\times}$.}

\vspace{45mm}

\noindent {\ubuntub (2.)} {\ubuntu Interestingly, $(\mathbb{Z}/p\mathbb{Z})^{\times}$ is always cyclic if $p$ is a prime number. Thus, you can find a generator for $(\mathbb{Z}/7\mathbb{Z})^{\times}$. What is it? There might be more than one.}

\end{document}
